\documentclass[12pt,a4paper]{article}

\usepackage{fontspec}
\usepackage{polyglossia}
\setmainlanguage{vietnamese}
\setotherlanguage{english}

\usepackage{geometry}
\geometry{margin=2.5cm}

\usepackage{hyperref}
\hypersetup{
  pdftitle={DAIOF Framework - Alpha Omega},
  pdfauthor={Nguyễn Đức Cường (alpha\_prime\_omega)},
  colorlinks=true,
  linkcolor=blue,
  urlcolor=blue
}

\usepackage{titling}
\usepackage{parskip}

\setmainfont{DejaVu Serif}
\setsansfont{DejaVu Sans}
\setmonofont{DejaVu Sans Mono}

\title{\textbf{DAIOF Framework}\\[0.5em]
       \large Alpha Omega — Tài Liệu Kỹ Thuật}
\author{Nguyễn Đức Cường (alpha\_prime\_omega)}
\date{\today}

\begin{document}

\maketitle
\tableofcontents
\newpage

\section{Giới Thiệu}

\textbf{DAIOF} (Digital AI Organism Framework) là một hệ sinh thái AI tự trị được thiết kế
để vận hành theo nguyên tắc bốn trụ cột cốt lõi:

\begin{enumerate}
  \item \textbf{An toàn (Safety)} — mọi hành động đều có kế hoạch rollback.
  \item \textbf{Đường dài (Long-term)} — xây dựng hệ thống bền vững.
  \item \textbf{Tin số liệu (Data-driven)} — quyết định dựa trên bằng chứng.
  \item \textbf{Hạn chế rủi ro (Risk Management)} — nhận diện rủi ro sớm.
\end{enumerate}

\section{Kiến Trúc Hệ Thống}

DAIOF được tổ chức theo mô hình phân tầng:

\subsection{Tầng Chính Sách (Policy Layer)}
Định nghĩa các quy tắc bất biến và nguyên tắc hoạt động của toàn hệ sinh thái.

\subsection{Tầng Quản Trị (Governance Layer)}
Điều phối các quyết định chiến lược và phê duyệt thay đổi quan trọng.

\subsection{Tầng Vận Hành (Runtime Layer)}
Thực thi các tác vụ tự động, giám sát trạng thái và ghi nhận bằng chứng.

\subsection{Tầng Hạ Tầng (Infrastructure Layer)}
Cung cấp tài nguyên tính toán, lưu trữ và mạng cho toàn hệ thống.

\section{Giao Thức OSLF}

Giao thức \textbf{OSLF} (Observe, Score, Learn, Finalize) được kích hoạt tự động
khi độ tin cậy của quyết định dưới 80\%.

\subsection{Giai Đoạn A — Phân Tích}
Phân tách yêu cầu, liệt kê giả định và kiểm tra an toàn.

\subsection{Giai Đoạn B — Tập Trung}
Chấm điểm từng phương án theo Bốn Trụ Cột (0--10 điểm mỗi trụ).

\subsection{Giai Đoạn C — Kiến Trúc Lại}
Tạo đề xuất tối ưu, tính điểm rủi ro (0--5) và đưa ra quyết định.

\section{Thuộc Tính}

Framework được tạo bởi \textbf{Nguyễn Đức Cường} (alpha\_prime\_omega).\\
Ngày tạo gốc: 30 tháng 10 năm 2025.\\
Phiên bản: 1.0.0.\\
Giấy phép: MIT.

\bigskip
\begin{center}
  \textit{Powered by HYPERAI Framework}\\
  \textit{Creator: Nguyễn Đức Cường (alpha\_prime\_omega)}
\end{center}

\end{document}
